\section{Introduction}
\label{sec:intro}

\looseness -1
Diffusion models are the dominant generative model class for
high-dimensional continuous data such as images, video, audio, and
molecules~\citep{ho2020denoising, song2021scoresde,
rombach2022highresolution, karras2022elucidating,
leeGenMolDrugDiscovery2025}. Sampling corresponds to numerically
integrating the time-reversed stochastic differential equation
associated with a forward noising process; the only learnable
component is a score network
$s_\theta(x_t, t) \approx \nabla \log p_t(x_t)$ trained to match the
noisy marginals of the data distribution. Sampling can equivalently
be expressed as a deterministic ordinary differential equation with
the score as a velocity field, an approach taken by flow-matching
methods~\citep{lipman2023flow}.

\looseness -1
Latent diffusion models~\citep{rombach2022highresolution, esser2024sd3,
podell2023sdxl} run forward and backward diffusion not in pixel space
but in a learned latent space of much smaller dimensionality $\dlat$,
typically obtained from a VAE encoder. By placing diffusion in a
representation closer to the data's true intrinsic dimensionality
$\dint$~\citep{pope2021, brown2022manifold, stanczuk2024manifold,
kamkari2024geometric}, though still well above it (with
$\dlat \gg \dint$ in practice), latent diffusion achieves comparable
or better sample quality at a fraction of the compute of pixel-space
diffusion.

\looseness -1
A model achieving zero training loss without regularization learns
the empirical score and reproduces training samples at the end of
the backward process~\citep{gu2023memorization}; the regime kicks
in when the training set is small relative to model capacity and
disappears past a dataset-size threshold~\citep{kadkhodaie2023learning}. What is missing is
a sharp account of how this transition depends on the latent-space
dimensionality, the design knob that distinguishes
latent-diffusion from pixel-space diffusion. We ask:
\emph{how does the gap between $\dlat$ and $\dint$ shape the
transition from generalization to memorization, and is there a
mechanism a practitioner can tune by adjusting the latent width?}

\looseness -1
Existing theory provides a starting point.
\citet{bonnaire2025} characterize diffusion training as a
competition between a generalization time $\taugen$, at which the
score network learns the population distribution, and a
memorization time $\taumem$, at which it begins fitting individual
training points. Their analysis assumes isotropic data covariance,
which collapses the distinction between the
$\dint$-dimensional signal subspace and the $\dlat - \dint$ null
directions of the latent space. We extend their spectral analysis
to the latent-diffusion regime $\dlat > \dint$, identify a new
noise-dimension bulk that opens between the population and sample
bulks, and show that it acts as a buffer that delays memorization.

\textbf{Contributions.}
\begin{itemize}\itemsep1pt
  \item We extend Bonnaire's two-bulk theorem to anisotropic data
    and identify a four-bulk eigenvalue structure for the score
    network's feature correlation matrix, with bulk-count
    boundaries that land exactly at indices $\dint$ and $\dlat$.
  \item We identify the noise-dimension bulk as a memorization
    buffer: increasing $\dlat$ adds $\dlat - \dint$ extra modes
    that the readout must absorb before reaching the sample bulk,
    delaying $\taumem$ while leaving $\taugen$ approximately
    fixed.
  \item We confirm the buffer mechanism in trainable MLP score
    networks on synthetic data with five-seed confidence intervals,
    holding hidden width fixed so the $\dlat$ sweep is not a model
    capacity sweep.
  \item We test the same idea in VAE latent spaces on CelebA and
    CIFAR-10. Memorization first rises as the VAE becomes faithful,
    then decreases at high $\dlat$, while FID is best at
    intermediate latent widths.
  \item We introduce a data-only spectral predictor: compute the
    VAE latent covariance spectrum, calibrate a single global
    pressure curve, and overlay predicted memorization thresholds on
    empirical memorization trajectories.
\end{itemize}
